% Promoted after complete root and reduction_descent proof/source audits; candidate preserved.
\subsection{Removing a common two-power from the polarization}

\begin{lemma}[The actual maximal two-power root]
\label{lem:DE-A3-mixed-common-power-root}
\status{Proved; complete independent proof and original-source audits passed.}
Let $A$ be an abelian surface of $2$-rank one over an algebraically
closed field of characteristic two, and let $M$ be ample.  Use the
integral polarization indices $a=a(M)$ and $b=b(M)$ of
Lemma~\ref{lem:DE-A3-mixed-polarization-blocks}, and put
\[
             r=\min\{a,b\},\qquad n=2^r,\qquad D=\chi(M).
\]
There is an ample line bundle $H$ on the same surface such that
\[
 M\simeq H^n,\qquad
 a(H)=a-r,\qquad b(H)=b-r,\qquad \chi(H)=D/n^2.
                                                        \tag{DE-CP1}
\]
Moreover $r$ is the largest nonnegative integer for which $M$ has
an actual $2^r$th root in $\Pic(A)$.  The root is not asserted to be
defined over a smaller, nonalgebraically closed ground field.
\end{lemma}
\begin{proof}
We use the actual integral blocks, not merely the rational slope
decomposition.  On each of the ordinary and local-local Dieudonne
summands, the matrix of $\lambda_M$ has both elementary divisors
divisible by $n$.  Equivalently it annihilates the entire finite
group scheme $A[n]$.  The integral splitting and exact Dieudonne
dictionary proving this equivalence were established in
Lemma~\ref{lem:DE-A3-mixed-polarization-blocks}; all its hypotheses
hold here.

The finite faithfully flat quotient $[n]:A\to A$ is the quotient
by $A[n]$.  Consequently there is a homomorphism
$\varphi:A\to A^\vee$ with
\[
                   \lambda_M=\varphi\circ[n]=n\varphi.
\]
Dualizing gives $n\varphi^\vee=n\varphi$.
The group $\Hom(A,A^\vee)$ has no integer torsion: if $n f=0$,
the morphism $f$ factors through the finite group scheme
$A^\vee[n]$, and a morphism from the proper geometrically
connected reduced variety $A$ to a finite affine scheme is constant;
as a group homomorphism it is zero.  Thus $\varphi$ is symmetric.

Every symmetric homomorphism of an abelian variety over an
algebraically closed field is induced by a line bundle, in arbitrary
characteristic, including two.  We use the full proof of
\cite[Proposition~11.2 and Corollary~11.3, pp.~159--160]
{EdixhovenVanDerGeerMoonenAV}.  That proof obtains the line bundle
by finite-flat isotropic descent from a Poincare pullback and does
not divide a theta-group argument by two.
Choose $H_0$ with $\lambda_{H_0}=\varphi$.
Then $P=M\otimes H_0^{-n}$ lies in $\Pic^0(A)$, since its associated
homomorphism is zero.  Multiplication by $n$ on the dual abelian
variety is surjective, hence surjective on algebraically closed
points.  There is therefore $Q\in\Pic^0(A)$ with $Q^n\simeq P$.
Set $H=H_0\otimes Q$.  This gives the actual isomorphism
$H^n\simeq M$.  A line bundle is ample if a positive tensor power
is ample, so $H$ is ample.

The equality of polarizations $\lambda_M=n\lambda_H$ subtracts
$r$ from each integral elementary-divisor valuation, proving both
index formulas.  Riemann--Roch on a surface gives
$\chi(H^n)=n^2\chi(H)$, proving the Euler-characteristic formula.
Conversely, if $M=H'^{\,2^s}$ for a line bundle $H'$, it is ample
and its integral polarization blocks have nonnegative indices.
The equality $\lambda_M=2^s\lambda_{H'}$ then gives $a,b\ge s$.
Thus $s\le r$, proving maximality.
\end{proof}

\begin{theorem}[All degrees when the ordinary index is no larger]
\label{thm:DE-A3-mixed-index-ordered-ranks}
\status{Proved; complete independent proof and original-source audits passed.}
Let $A$ be an abelian surface of $2$-rank one over a perfect field
of characteristic two, and let $M$ be ample.  Compute $a(M),b(M)$
after algebraic closure, and assume
\[
                         0<a(M)\le b(M).
\]
Put $D=\chi(M)$ and $L=M^3$.  For every $d\ge1$, the geometric
generic relative stable rank in the complete system $|L^d|$ is
\[
                             r_d=9Dd^2-2.              \tag{DE-CP2}
\]
The relative coherent dimension is $9Dd^2-1$, and the generic
smooth curve has genus $9Dd^2+1$ and $2$-rank $9Dd^2-1$.
For the embedding by the very ample bundle $L$, the same relative
stable rank holds for ambient degree-$d$ hypersurfaces whenever
restriction to the complete system is surjective; in particular
the earlier bound
$d\ge\max\{1,\operatorname{reg}(\mathcal I_A)-1\}$ suffices.
\end{theorem}
\begin{proof}
Work over an algebraic closure, as permitted by the geometric
generic assertion.  Apply the preceding lemma with $r=a(M)>0$.
It gives an actual ample root $H$ on $A$ with $a(H)=0$ and
$M\simeq H^{2^a}$.  The complete systems are therefore identical:
\[
             |M^{3d}|=|H^{3d'}|,\qquad d'=2^a d.
\]
The integer $d'$ is even.  The full proof of
Theorem~\ref{thm:DE-A3-mixed-nonproduct-ranks} applies to $(A,H)$:
it supplies sharp odd-degree nodal norms from an etale principal
source, obtains every even degree by transverse unions, and uses
the integral Prym obstruction for the matching upper bound.
In particular its even formula gives
\[
 r_d=9\chi(H)(d')^2-2
     =9(D/2^{2a})(2^a d)^2-2
     =9Dd^2-2.
\]
All systems here are the full systems of the same actual line
bundle; no claim of genericity under a restricted pullback map is
needed.  Adjunction, ample vanishing, and subtraction of the ambient
stable rank one give the remaining dimensions, as in that theorem.
The bundle $M^3$ is very ample.  Surjectivity from ambient
hypersurfaces, with geometrically integral fibres on coefficient
spaces, transfers the same generic rank, exactly as in the already
proved regularity argument.
\end{proof}

\begin{corollary}[A primitive form of the remaining mixed problem]
\label{cor:DE-A3-mixed-primitive-rank-gate}
\status{Proved exact reduction; primitive-case nonvanishing remains open.}
For every ample $M$ on a mixed abelian surface, the same-surface
root $H$ of (DE-CP1) identifies the degree-$d$ rank problem for
$M^3$ with the degree-$2^r d$ problem for $H^3$.
Thus, together with the proved ordinary-index-zero theorem and
(DE-CP2), the unresolved nonproduct positive-index case reduces to
ample bundles satisfying
\[
                         a(H)>0,\qquad b(H)=0,
                                                        \tag{DE-CP3}
\]
in all positive degrees.  This reduction does not prove the
restricted cup-product radical is zero, and does not assert that
an arbitrary mixed surface has an index-zero ample class.
\end{corollary}
\begin{proof}
The actual line-bundle equality identifies the systems and their
families, hence their geometric generic Frobenius ranks.  If
$a(M)\le b(M)$, its root has ordinary index zero and the relevant
degree has the parity already used above.  If $a(M)>b(M)$, take
$r=b(M)$.  Its root has the indices in (DE-CP3), and the exponent
changes by the displayed factor.  The root stays on the original
surface; only its line bundle and the exponent change.
\end{proof}
